
\documentclass[11pt]{scrartcl}

\usepackage[margin=1in]{geometry}
\usepackage{setspace}
\usepackage{booktabs}
\usepackage{longtable}
\usepackage{enumitem}
\usepackage{listings}
\usepackage{xcolor}
\usepackage{tikz}
\usetikzlibrary{positioning, arrows.meta, shapes.multipart, fit}

\usepackage[colorlinks=true,
            linkcolor=black,
            urlcolor=blue,
            citecolor=black]{hyperref}

\usepackage{amsmath,amssymb}

\setstretch{1.0}
\setlength{\parskip}{0.3em}
\setlength{\parindent}{0pt}

% --- KOMA-native (do NOT use titlesec with scrartcl) ---
\RedeclareSectionCommand[beforeskip=1.0ex plus .2ex minus .2ex, afterskip=0.4ex plus .1ex]{section}
\RedeclareSectionCommand[beforeskip=0.8ex plus .2ex minus .2ex, afterskip=0.3ex plus .1ex]{subsection}
\RedeclareSectionCommand[beforeskip=0.6ex plus .2ex minus .2ex, afterskip=0.2ex plus .1ex]{subsubsection}

% ---------- listings ----------
\definecolor{codegray}{gray}{0.96}
\definecolor{coderule}{gray}{0.65}
\lstset{
  backgroundcolor=\color{codegray},
  basicstyle=\ttfamily\small,
  breaklines=true,
  frame=single,
  rulecolor=\color{coderule},
  columns=fullflexible,
  showstringspaces=false
}

\title{Architectural Refactor Technical Report \\ FUSE Modeling Framework}
\author{Martyn Clark (refactor summary draft)}
\date{\today}

\begin{document}

\maketitle
\tableofcontents

% =====================================================
\section*{Abstract}

\begin{abstract}

This report documents a major capability expansion of the FUSE hydrologic modeling
framework. The new release introduces a fully differentiable physics pathway with
analytical Jacobians throughout the snow, soil, and routing components; a purpose-built
fixed-step implicit solver designed for deterministic, derivative-consistent execution;
and a generalized spatial-domain abstraction supporting lumped, distributed,
rectilinear, and curvilinear configurations. The framework now also supports structured
spatial decomposition, enabling scalable parallel execution.

Implementing these capabilities required a substantial architectural re-platforming.
The refactor spans 95 source files and replaces implicit global state with explicit,
typed run containers (\texttt{info}, \texttt{work}, \texttt{domain}), modular
driver/evaluation logic, explicit time-axis construction, improved NetCDF I/O and
domain handling, and modernized calibration plumbing. A temporary compatibility
layer preserves selected legacy interfaces during migration but is not part of the
target architecture.

Beyond describing \emph{what} changed, this report explains \emph{why} the legacy
structure became a limiting factor under differentiable and scalable execution
requirements, and \emph{how} the redesigned architecture improves correctness,
traceability, maintainability, and extensibility for future scientific development.

\end{abstract}
% =====================================================
\section{Background and Motivation}

\subsection{From Refactor to Capability Expansion}

This development effort was not initiated solely to ``clean up'' legacy code.
It was driven by the need to support a set of substantial new capabilities:

\begin{itemize}
\item A fully differentiable physics pathway with analytical Jacobians.
\item A purpose-built fixed-step implicit solver suitable for gradient-based workflows.
\item Flexible spatial-domain support (lumped, distributed, rectilinear and curvilinear grids).
\item Structured spatial decomposition enabling parallel execution.
\end{itemize}

These capabilities fundamentally expand what FUSE can do: from a forward
simulation framework with flexible numerical options to a platform capable of
derivative-aware modeling, scalable spatial execution, and integration with
modern calibration and machine-learning workflows.

\subsection{Architectural Shift: From Implicit Globals to Explicit State}

Implementing these capabilities required more than incremental feature additions.
They imposed stricter requirements on determinism, state management, numerical
stability, and data flow transparency. Differentiable execution and spatial
decomposition, in particular, require predictable control flow and well-defined
state ownership across the entire call graph.

In particular, correctness often depended on initialization order and
side effects rather than explicit interfaces. This created challenges for:

\begin{itemize}
\item Isolating and testing components,
\item Safely modifying numerical routines,
\item Supporting deterministic execution pathways,
\item Enabling derivative-consistent integration.
\end{itemize}

The central architectural change is therefore a shift in programming model:
from code that operates through implicit shared state to code that
propagates explicitly defined, typed state containers through well-defined
interfaces.

This change is not aesthetic. It provides the structural foundation required
for differentiable physics, fixed-step implicit integration, and scalable
spatial execution.

\subsection{Guiding principles}

The redesign follows several principles commonly used in large scientific codes:

\begin{enumerate}[leftmargin=*]
\item \textbf{Make state explicit:} The state required to run the model should be visible
      in the subroutine interfaces (or in a clearly defined container passed in).
\item \textbf{Separate definition from storage:} Define data structures in one place
      (\texttt{types/}), and store/instantiate them in clearly owned containers
      (\texttt{info}, \texttt{work}, \texttt{domain}).
\item \textbf{Keep orchestration separate from execution:} The driver should orchestrate
      setup and mode selection; evaluation should execute the model and compute metrics.
\item \textbf{Enable incremental migration:} Preserve legacy call sites temporarily so
      the refactor can be staged; remove the compatibility layer once migration is complete.
\end{enumerate}

\section{Development of New Capabilities}

\subsection{Differentiable Physics Architecture}

The differentiable branch is not a minor extension of the legacy physics;
it is a systematic reworking of the process layer to support derivative-aware
simulation and gradient-based workflows.

The refactor introduces a parallel set of \texttt{*\_diff} routines that
propagate analytical derivatives throughout the snow, soil, and routing
physics. Rather than relying on external automatic differentiation or
finite-difference approximations, derivatives are computed explicitly
within the governing equations.

\subsubsection{Fully Differentiable Snow Model}

The differentiable snow implementation (\texttt{update\_swe\_diff}) is a derivative-aware
rewrite of the elevation-band SWE update.
The routine computes snowfall, rainfall, melt, and SWE evolution for each elevation band,
while propagating analytical derivatives with respect to the snow-related parameters.

For each band, forcing is adjusted using (i) a temperature lapse rate and (ii) an orographic
precipitation gradient (OPG) multiplier based on the band--forcing elevation difference.
Precipitation is then partitioned into snow and rain using a smooth sigmoid transition
around the temperature threshold \texttt{PXTEMP}, yielding a differentiable snow fraction
\(f_{\text{snow}}\in[0,1]\). Melt is computed using a seasonally varying melt factor
\(MF\) derived from \texttt{MFMIN} and \texttt{MFMAX} via an annual sinusoid (with explicit
leap-year handling), multiplied by a \emph{smoothed} positive-temperature term
\(\max(T - \texttt{MBASE},0)\) implemented with \texttt{smax} to ensure stable derivatives.

Two deliberate non-smooth elements are retained for physical realism and numerical hygiene:
(i) melt is hard-capped by available SWE (\(\texttt{snowmelt}=\min(\texttt{potMelt},\texttt{SWE}/\Delta t)\))
to avoid negative SWE, and (ii) SWE is hard-clamped at zero to remove small numerical noise
and prevent ``ghost snow''. Sensitivities are correspondingly gated to zero during snow-free
periods, avoiding accumulation of meaningless derivatives when SWE is effectively absent.

A key capability addition is the consistent propagation of parameter sensitivities through
the entire banded update. The routine carries forward stored sensitivities
(\texttt{dSWE\_dparam}) from the previous timestep, updates them using the analytical
derivatives of accumulation and melt, and aggregates the resulting derivatives into the
band-area-weighted effective precipitation flux (\(\texttt{EFF\_PPT} = \texttt{rain} + \texttt{snowmelt}\)).
This provides a differentiable forcing-to-effective-precipitation mapping suitable for
gradient-based calibration and coupled differentiable modeling workflows.


\subsubsection{Analytical Jacobians in Soil Physics}

The soil moisture and flux equations were restructured to expose and
compute analytical partial derivatives with respect to state variables
and parameters. This includes:

\begin{itemize}
\item Explicit Jacobian terms for infiltration, percolation, interflow,
      baseflow, and saturation-excess runoff.
\item Derivative-aware implementations of constitutive relationships
      (e.g., soil moisture–conductivity functions).
\item Consistent propagation of sensitivities through state updates.
\end{itemize}

These Jacobians are integrated directly into the time-stepping logic,
enabling stable sensitivity propagation and forming the basis for
adjoint or gradient-based calibration frameworks.


\subsubsection{Architectural Significance}

This differentiable pathway transforms the physics layer from a purely
forward simulator into a gradient-aware dynamical system. It enables:

\begin{itemize}
\item Gradient-based parameter estimation.
\item Efficient high-dimensional calibration.
\item Integration with adjoint or automatic differentiation frameworks.
\item Future coupling to machine-learning components.
\end{itemize}

The scale of this change is substantial: rather than wrapping legacy code
with differentiation tools, the physics equations themselves were rewritten
to expose analytical structure. This represents a foundational capability
expansion, not merely a refactor.

\subsection{Fixed-Step Implicit Solver for Differentiable Execution}

FUSE has long included a broad suite of numerical options (explicit and implicit schemes,
adaptive sub-stepping, and multiple solution strategies; Clark and Kavetski, WRR, 2010).
The differentiable branch, however, places different requirements on the integration
layer: deterministic fixed time stepping (to avoid step-size branching), stable nonlinear
solves, and derivative-consistent state trajectories.

To meet these requirements, the refactor adds a dedicated fixed-step implicit solver
(\texttt{implicit\_solve\_module}). 

\subsubsection{Newton-Raphson solver with constraint handling}

The solver advances the state using an implicit Euler
update written as a nonlinear root-finding problem:
\[
F(x) \equiv x - \left(x_0 + \Delta t\, g(x)\right) = 0,
\]
where \(g(x) = \mathrm{d}x/\mathrm{d}t\) is provided by the differentiable physics layer.

The solver computes an analytical flux Jacobian \(J_g(x)\) directly from the differentiable
physics routine (\texttt{mod\_derivs\_diff}), and forms the residual Jacobian
\[
J(x) = \frac{\partial F}{\partial x} = I - \Delta t\, J_g(x).
\]
Each Newton iteration then solves \(J\,\delta x = -F\) using an LU factorization
(\texttt{ludcmp}/\texttt{lubksb}). This design avoids finite-difference Jacobians during the
core solve (a finite-difference option exists as a diagnostic fallback but is not the
default pathway).

A central feature of the implementation is robust handling of bounded state variables.
The solver explicitly computes state bounds (\texttt{get\_bounds}), applies soft
overshoot correction (\texttt{fix\_ovshoot}) when needed, and enforces feasibility during
step acceptance. In particular:

\begin{itemize}
\item Candidate updates that violate bounds are rejected during line search, and the step
      length is shrunk until the trial state lies in the feasible region.
\item An active-set style modification prevents Newton steps from pushing already-active
      bound constraints further outward.
\end{itemize}

This ``feasibility-first'' approach is critical for hydrologic states with physical bounds
(e.g., storages) and is especially important for differentiable execution, where repeated
bound violations can produce unstable trajectories and unusable gradients.

\subsubsection{Backtracking line search and multi-level fallback strategies.}
Rather than relying on an unconstrained Newton step, the solver uses a short backtracking
line search to reduce the implicit residual norm
\(\phi(x) = \tfrac{1}{2}\,F(x)^{\mathsf{T}}F(x)\), while maintaining feasibility. The solver
also implements several explicit fallback mechanisms to improve robustness:

\begin{itemize}
\item Step-size limiting (\(\|\delta x\|\) capped by a maximum step size).
\item If the Newton direction is not a descent direction for \(\phi\), a fallback direction
      based on the residual-gradient (\(-\nabla \phi\)) is used.
\item If line search fails to find an acceptable step, a damped trial step is taken and
      overshoot is repaired via soft clamping.
\item The solver tracks the best (lowest-\(\phi\)) feasible state encountered. In the rare
      case of non-convergence, it falls back to the best feasible evaluation, and (if
      necessary) uses an explicit Euler update with mass-conserving clamp/disaggregation
      (\texttt{conserve\_clamp}) to ensure physically consistent states.
\end{itemize}

\subsubsection{Significance.}
This solver is intentionally specialized: it targets fixed-step, derivative-consistent,
constraint-aware integration suitable for differentiable physics and gradient-based
workflows. It complements (rather than replaces) the flexible ``kitchen sink'' numerical
framework developed for forward simulation in FUSE, while providing a robust and
deterministic integration pathway for differentiable execution.

\subsection{Unified Spatial Domain Abstraction}

The refactor introduces a generalized spatial-domain abstraction that supports:

\begin{itemize}
\item Lumped catchments (single spatial point)
\item Distributed catchments (arbitrary point lists / HRU collections)
\item Gridded domains (rectilinear and curvilinear grids)
\end{itemize}

Previously, spatial assumptions were embedded implicitly throughout the codebase.
Many routines assumed specific dimensional layouts, and extending from lumped to
distributed configurations required invasive changes.

The new \texttt{domain} container centralizes spatial metadata (dimensions,
indexing, and decomposition logic), allowing evaluation and physics routines
to operate generically over spatial elements. This eliminates hard-coded
assumptions about domain shape and makes spatial flexibility a structural feature
rather than a patchwork of conditional logic.

\subsection{Parallelization via Spatial Decomposition}

The refactor also establishes a clear pathway for parallel execution by
decomposing the second spatial dimension (e.g., via \texttt{ystart}-style
index partitioning).

Rather than assuming a monolithic spatial loop, the domain can now be
partitioned into subdomains, enabling:

\begin{itemize}
\item Parallel execution across spatial slices
\item Improved memory locality
\item Scalable execution for large grids
\end{itemize}

Crucially, this decomposition is coordinated through the centralized
\texttt{domain} abstraction, rather than being embedded implicitly
in physics or I/O routines. This makes parallelism an architectural
extension of the model rather than an intrusive modification.

Together, these changes transform spatial handling from a fixed structural
constraint into a flexible, scalable capability.

\subsection{Structural improvements}

The refactoring includes:

\begin{enumerate}

  \item Explicit Run-State Architecture (\texttt{info/work/domain})

The transition from implicit module-level globals to explicitly passed run containers
(\texttt{info}, \texttt{work}, \texttt{domain}) is the most consequential improvement.
In the legacy architecture, routines depended on shared modules, and true dependencies
were not visible in argument lists. Correctness relied on initialization order and
side effects.

The refactor makes state ownership explicit and routes execution through a clean
evaluation interface:

\begin{lstlisting}[language=Fortran]
call fuse_evaluate(XPAR, info, work, domain, OUTPUT_FLAG, METRIC_VAL)
\end{lstlisting}

This improves traceability, testability, and reuse across deterministic and calibration modes.

\item Configuration Modernization: \texttt{fuseFileManager} $\rightarrow$ TOML

The legacy configuration pathway accumulated parsing, defaults, filesystem logic,
and runtime behavior inside a large monolithic routine. Configuration was effectively
encoded in control flow.

Transitioning to TOML introduces a structured, human-readable configuration layer
with earlier validation and clearer separation between configuration and execution.
This substantially improves reproducibility and maintainability.

\item NetCDF and Dimensional Cleanup

The refactor also addresses structural problems in forcing and output routines,
including monolithic I/O logic and hard-coded repetition (e.g., elevation-band writes).
Moving toward dimension-aware abstractions and smaller helper routines reduces
duplication and restores logical structure to I/O subsystems.

Two related cleanups substantially reduce legacy complexity. First, the refactor deletes
\texttt{get\_modtim} and converts the old ``time I/O'' module into a pure calendar utility
module: \texttt{time\_io.f90} becomes \texttt{util/time\_utils.f90}, drops the \texttt{netcdf}
dependency, and exports only \texttt{date\_extractor}, \texttt{juldayss}, and \texttt{caldatss}.
This removes slow, repetitive per-timestep NetCDF access and aligns time handling with the
new pattern: read the full time axis once, build \texttt{info\%time\%jdate(:)} once, and
convert any timestep via \texttt{caldatss(info\%time\%jdate(itim),\ldots)}.

Second, the refactor removes the legacy forcing-descriptor pathway (e.g., deleting
\texttt{force\_info.f90} and the versioned \texttt{forcinginfo} text formats), which had grown
into a large body of fragile code whose primary job was to parse \emph{descriptors of
self-describing NetCDF files}. With TOML-based configuration, any true overrides/aliases can
live in TOML, while the NetCDF metadata remains the authoritative source for forcing
definitions. Together, these changes eliminate substantial ``meta-I/O'' code, reduce hidden
coupling, and make initialization faster and easier to reason about.

\end{enumerate}

These changes reduce ``incidental complexity'' in the codebase while enabling new
capabilities (clean configuration, cleaner runtime diagnostics, and differentiable execution).

This shift from ad hoc procedural growth toward structured, dimension-aware
design reduces duplication, clarifies intent, and restores logical coherence
to core subsystems such as forcing ingestion and output generation.


% =====================================================
\section{Review of the Legacy Architecture}

\subsection{FUSE Architecture: What it was and why it became suboptimal}

The legacy FUSE organization relied heavily on shared modules (e.g., \texttt{dshare/*})
as global repositories for forcing data, parameters, states, routing, and statistics.
Many routines relied on \texttt{USE multiforce}, \texttt{USE multistate}, etc., and
implicitly expected those globals to be initialized before use.

This style has advantages for small/medium codes:
\begin{itemize}[leftmargin=*]
\item It reduces argument lists and avoids passing many arrays around.
\item It allows quick prototyping when many routines need access to shared state.
\end{itemize}

However, as the codebase grows and more run modes appear, the disadvantages dominate.

\subsection{Why global-module state is suboptimal at scale}

\subsubsection{Hidden dependencies and side effects}

In a global-state design, a subroutine’s argument list does not fully describe the
state it reads or modifies. For example:

\begin{lstlisting}[language=Fortran]
call get_modtim(ncid, itim, ierr, message)
\end{lstlisting}

Although this call appears to simply compute a timestamp, it may also modify
module-level variables (e.g., \texttt{timDat\%dtime}, \texttt{timDat\%iy},
or \texttt{jdayRef}). These variables are then implicitly used by other
routines --- for example, within a snow model or energy-balance calculation ---
without any indication in their argument lists that the values were updated.

As a result, downstream routines depend on side effects that are not visible
at the call site. If \texttt{get\_modtim} is removed, reordered, or replaced
during refactoring, model behavior can change even though the replacement code
appears functionally equivalent. This hidden coupling creates a common failure
mode: \emph{A small, local change produces a large behavioral difference, and
the rootcause is not obviously connected to the modified lines}.

\subsubsection{Testing and reproducibility are harder}

Global-state code makes unit tests difficult because isolating a routine requires
constructing a large, often implicit environment (multiple global modules initialized
in the correct order). Regression tests also become harder to interpret: if a change
affects results, it can be unclear whether the difference is due to the intended change
or an unintended shift in global state usage.

\subsubsection{Extensibility costs grow nonlinearly}

Adding new features is harder when:
\begin{itemize}[leftmargin=*]
\item A new forcing format requires changes across many routines that read globals.
\item Domain decomposition changes how arrays are shaped; global arrays become ambiguous
      (global vs local indexing).
\item New evaluation logic or calibration schemes require threading new metadata through
      a chain of routines that do not accept explicit arguments.
\end{itemize}

At some point the main cost is no longer implementing the new feature, but ensuring it
does not break unrelated functionality.

\subsection{Structural Growth and Maintainability Issues in the Legacy Architecture}

Beyond global state coupling, the legacy codebase exhibited a pattern of
organic growth without consistent structural boundaries. Over time, large
routines accumulated functionality well beyond their original intent,
often combining multiple responsibilities without clear logical separation.

\subsubsection{Monolithic Routines and Task Accretion}

Several legacy routines expanded to perform a wide range of loosely related
tasks. For example, the previous \texttt{get\_gforce} routine handled:

\begin{itemize}
\item NetCDF variable lookup
\item Dimension queries
\item Domain subsetting
\item Elevation band handling
\item Unit interpretation
\item Forcing extraction and reshaping
\end{itemize}

These operations were implemented in a largely sequential, ad hoc manner,
with minimal modular decomposition. As a result:

\begin{itemize}
\item Logical structure was difficult to discern.
\item Debugging required tracing through hundreds of lines of mixed concerns.
\item Small changes (e.g., adding a new forcing variable) required edits
      in multiple scattered locations.
\item Code readability degraded as the routine expanded.
\end{itemize}

The routine effectively became a ``catch-all'' forcing processor rather than
a composable set of focused operations.

\subsubsection{Hard-Coded Repetition Instead of Dimensional Abstraction}

A more troubling structural issue appeared in output routines such as
\texttt{def\_output} and \texttt{put\_output}. Instead of leveraging
NetCDF dimensions for elevation bands, the legacy implementation contained
dozens of nearly identical write statements, one per band. For example,
rather than defining:

\begin{itemize}
\item a dimension \texttt{elevband}
\item and writing arrays indexed by that dimension,
\end{itemize}

the code manually wrote each elevation band variable individually,
resulting in patterns such as:

\begin{itemize}
\item repeated \texttt{nf90\_put\_var} calls for band 1
\item repeated \texttt{nf90\_put\_var} calls for band 2
\item \ldots
\item repeated \texttt{nf90\_put\_var} calls for band 50
\end{itemize}

This design introduces several serious drawbacks:

\begin{enumerate}
\item \textbf{Scalability limitations:} Increasing the number of elevation bands
      requires adding new code blocks rather than simply changing a dimension.
\item \textbf{High maintenance cost:} Any modification to output logic must be
      replicated consistently across many nearly identical sections.
\item \textbf{Increased error risk:} Copy–paste repetition increases the
      probability of subtle inconsistencies.
\item \textbf{Lost abstraction opportunity:} NetCDF is dimension-oriented by design;
      bypassing this structure undermines clarity and efficiency.
\end{enumerate}

\subsection{Consequences of Organic Growth}

As the codebase grew, these patterns accumulated:

\begin{itemize}
\item Routines performing too many responsibilities.
\item Hard-coded repetition where loops and dimensions should be used.
\item Logic arranged historically rather than conceptually.
\item Difficulty reasoning about control flow and data ownership.
\end{itemize}

The cumulative effect was not merely aesthetic; it increased the risk that
future extensions would further entrench complexity.


% =====================================================
\section{Target Architecture: What it is and why it is better}

The refactor introduces an explicit run-state model based on three main containers:

\begin{itemize}[leftmargin=*]
\item \textbf{\texttt{info}}: configuration and metadata (time axis, windows, units, CLI options,
      file IDs, model definition metadata).
\item \textbf{\texttt{work}}: runtime scratch and transient state (per-step buffers, intermediate
      arrays, solver workspace).
\item \textbf{\texttt{domain}}: spatial decomposition container (domain dimensions, mapping, local/global indices).
\end{itemize}

Instead of relying on global modules, routines now operate on these containers (or
on well-scoped subsets of them). This has immediate benefits:

\begin{enumerate}[leftmargin=*]
\item \textbf{Traceability:} data required by a routine is visible in its inputs.
\item \textbf{Correctness:} fewer implicit order dependencies; stale state is harder to accidentally use.
\item \textbf{Testability:} containers can be constructed for unit tests without initializing unrelated globals.
\item \textbf{Extensibility:} new fields can be added to types without rewriting every call site.
\end{enumerate}

\subsection{Separation of type definitions from storage}

The \texttt{types/} directory defines canonical structures:
\begin{itemize}[leftmargin=*]
\item \texttt{info\_types.f90}: defines \texttt{type(fuse\_info)} including \texttt{time\_info},
      \texttt{space\_info}, and run configuration.
\item \texttt{work\_types.f90}: defines runtime work arrays and scratch buffers.
\item \texttt{data\_types.f90}: defines domain-level storage (e.g., \texttt{domain\_data}).
\end{itemize}

This separation means that the meaning of fields is centralized, and multiple components
(driver, evaluation, I/O) can rely on the same type definitions.

Modules under \texttt{share/*\_data.f90} exist only to preserve legacy code paths
during incremental migration. They mimic previous module-level globals (e.g., \texttt{multiforce})
so that older routines continue to compile. They are intentionally not part of the end-state design.

The legacy architecture centered on implicit global state: many routines depended on
\texttt{dshare} globals and therefore on a correct initialization order. The target
architecture centralizes state in explicit containers passed through evaluation.

The compatibility layer remains temporarily to support incremental migration but is not
part of the end-state design.

% =====================================================
\subsection{Driver and Evaluation Refactor}

\subsubsection{Why the legacy driver structure was suboptimal}

Historically the driver mixed many concerns:
\begin{itemize}[leftmargin=*]
\item argument parsing / run mode selection
\item domain setup and model-definition setup
\item forcing ingestion and time indexing
\item model execution and statistics/metrics
\item calibration algorithm integration
\end{itemize}

When all of these are interleaved, changing one part can unexpectedly affect others.
For example, modifying evaluation logic might require touching the driver, which
also handles calibration and I/O side effects. This increases the risk that a change
introduces subtle differences (e.g., order of initialization, timing of exports to globals).

\subsubsection{Why the new split is better}

The refactor isolates responsibilities:
\begin{itemize}[leftmargin=*]
\item The \textbf{driver} orchestrates: it builds \texttt{info/domain/work}, selects a run mode,
      and calls evaluation or calibration.
\item The \textbf{evaluation} module executes: it runs the time loop, calls physics, and returns a metric.
\end{itemize}

This split has practical benefits:
\begin{enumerate}[leftmargin=*]
\item Evaluation becomes callable from multiple contexts (deterministic runs, SCE callback, future optimizers).
\item The driver becomes easier to read and reason about, because it does not contain the entire numerical workflow.
\item Testing becomes easier: evaluation can be tested independently of command-line parsing.
\end{enumerate}

\subsubsection{Illustrative interface}

The new evaluation entry point is conceptually:
\begin{lstlisting}[language=Fortran]
subroutine fuse_evaluate(XPAR, info, work, domain, OUTPUT_FLAG, METRIC_VAL)
\end{lstlisting}

This makes the dependencies explicit. In a container-based design, this signature also
provides a stable surface for new optimization tools and new run modes.

% =====================================================
\subsection{Time Axis and Windowing: From implicit state to explicit products}

\subsubsection{Why the legacy approach caused problems}

In legacy code, time handling was spread across multiple routines and modules. Some parts
derived time indices from configuration strings; other parts read NetCDF time variables;
still others computed calendar dates by mutating globals. The consequence is that the
\emph{time axis} (the sequence of times the model is actually run on) was not clearly
defined in one place.

This creates multiple failure modes:
\begin{itemize}[leftmargin=*]
\item Forcing time units (seconds/hours/days) may be interpreted inconsistently.
\item ``Window start/end'' logic can vary depending on how indices are mapped.
\item Removing a seemingly redundant time-reading call can change results because it also updated global structures.
\end{itemize}

\subsubsection{What the new code does}

The refactor treats time initialization as a deliberate pipeline:

\begin{enumerate}[leftmargin=*]
\item \textbf{Read} the NetCDF time coordinate and units attribute.
\item \textbf{Interpret} the units and compute a scale factor (seconds/minutes/hours $\rightarrow$ days).
\item \textbf{Build} the Julian-day axis in a single canonical representation.
\item \textbf{Map} simulation/evaluation date strings to indices in that axis.
\item \textbf{Validate} that evaluation is a subset of simulation and that windows fit within forcing coverage.
\item \textbf{Store} all outputs in \texttt{info\%time} as the single source of truth.
\end{enumerate}

\subsubsection{Illustrative code snippets}

\begin{lstlisting}[language=Fortran]
call read_time_axis(info%files%ncid_forc, info%time%time_steps, units_local, nt, ierr, cmessage)
call build_julian_axis(info%time%time_steps, trim(units_local), &
                       info%time%jdate_ref, info%time%jdate, info%time%deltim_days, ierr, cmessage)
\end{lstlisting}

\begin{lstlisting}[language=Fortran]
call map_dates_to_indices(info%time%jdate, date_start_sim, date_end_sim, &
                          info%time%sim_beg, info%time%sim_end, ierr, cmessage)
call map_dates_to_indices(info%time%jdate, date_start_eval, date_end_eval, &
                          info%time%eval_beg, info%time%eval_end, ierr, cmessage)
\end{lstlisting}

\subsubsection{Why this is better}

\begin{itemize}[leftmargin=*]
\item \textbf{Single source of truth:} there is one canonical time axis (\texttt{info\%time\%jdate}) and one
      canonical timestep (\texttt{info\%time\%deltim\_days}).
\item \textbf{Explicit unit handling:} scale factors are computed once, centrally, and documented.
\item \textbf{Window logic becomes auditable:} indices are computed in one place with validation checks.
\item \textbf{Fewer side effects:} time conversion routines are pure utilities rather than state-mutating I/O.
\end{itemize}

\subsection{Configuration Modernization: \texttt{fuseFileManager} $\rightarrow$ TOML}

\subsubsection{Why the legacy \texttt{fuseFileManager} approach was suboptimal}

The original configuration pathway centered on a large, monolithic file-management and
configuration routine. Over time, configuration logic accumulated in a way that made it
difficult to understand what was user-configurable versus what was derived at runtime.
This produced several recurring problems:

\begin{itemize}
\item \textbf{Poor separation of concerns:} configuration parsing, validation, defaulting,
      and filesystem logic were interleaved, increasing the cost of maintenance.
\item \textbf{Limited validation:} many configuration failures surfaced late (during I/O
      or execution) rather than early during parsing.
\item \textbf{Low readability:} the effective configuration schema was spread across code paths
      rather than being clearly represented as a structured document.
\end{itemize}

\subsubsection{Why TOML is an improvement}

Moving to TOML replaces the ``configuration-by-control-flow'' style with an explicit,
human-readable configuration document:

\begin{itemize}
\item \textbf{Explicit schema and intent:} settings live in a structured file rather than
      being encoded implicitly in procedural logic.
\item \textbf{Earlier failure modes:} parsing and validation can fail fast before expensive
      initialization occurs.
\item \textbf{Maintainability:} adding a new option becomes a small, localized change rather
      than a modification scattered across a large file-management routine.
\end{itemize}

This transition improves usability for end users and reduces coupling between configuration,
file-system conventions, and runtime model execution.



% =====================================================
\subsection{NetCDF I/O and Domain Decomposition}

\subsubsection{Why forcing I/O was a refactor hotspot}

Forcing ingestion sits at the intersection of multiple concerns:
\begin{itemize}[leftmargin=*]
\item NetCDF metadata (dims, variable IDs, units)
\item spatial mapping (HRUs, bands, cells)
\item time indexing and subsetting
\item memory layout and performance
\end{itemize}

In the legacy design, forcing logic often depended on global modules for dimensions and indices,
which makes it easy for mismatches to slip in when changing domain layout or adding decomposition.

\subsubsection{What the refactor improves}

The refactor introduces explicit helpers and domain-dimension utilities so that array shape and
dimension logic are concentrated and reusable. New files (e.g., \texttt{get\_domain\_dims.f90},
\texttt{domain\_decomp.f90}) are intended to formalize what ``the domain'' means in terms of indices,
sizes, and decomposition.

\subsubsection{Why this is better}

\begin{enumerate}[leftmargin=*]
\item \textbf{Reduced ambiguity:} ``global vs local'' dimensionality is explicit and queryable.
\item \textbf{Safer extensions:} new forcing variables can be added using shared lookup/query routines.
\item \textbf{Performance opportunities:} centralized dimension logic makes it easier to optimize reads and memory layout.
\end{enumerate}

% =====================================================
\subsection{Calibration and Optimization Plumbing}

\subsubsection{Why legacy optimization coupling was suboptimal}

Calibration code is typically ``outer loop'' logic that should not depend on the internal
details of forcing ingestion, time stepping, or global module state. In practice, the legacy
implementation relied on global modules, which makes it fragile:
\begin{itemize}[leftmargin=*]
\item The optimizer callback can accidentally depend on stale global state between evaluations.
\item Memory allocations in hot loops can be expensive and hard to notice.
\item It is difficult to reuse evaluation logic across different optimizers without duplicating code.
\end{itemize}

\subsubsection{Why a callback context is better}

The refactor introduces a callback context (\texttt{sce\_callback\_context}) that holds the pointers
to \texttt{info}, \texttt{work}, and \texttt{domain}. The SCE objective function can then call
\texttt{fuse\_evaluate} using that context. This improves:

\begin{itemize}[leftmargin=*]
\item \textbf{Clarity:} the callback clearly depends on the evaluation API, not on global state.
\item \textbf{Reusability:} other optimizers can use the same evaluation entry point.
\item \textbf{Performance:} parameter arrays and workspaces can be allocated once where appropriate.
\end{itemize}

\subsubsection{Illustrative snippet}

\begin{lstlisting}[language=Fortran]
! Pseudocode illustrating the intent of the callback context
ctx%info   => info
ctx%work   => work
ctx%domain => domain
call fuse_evaluate(par, ctx%info, ctx%work, ctx%domain, output_flag, metric_val)
\end{lstlisting}

% =====================================================
\subsection{Physics Layer Expansion and Differentiable Pathways}

\subsubsection{Why this was introduced}

Support for differentiable modeling and improved robustness often requires:
\begin{itemize}[leftmargin=*]
\item smoother or continuous approximations (to avoid discontinuities in derivatives)
\item stable implicit/iterative solves (for stiff or constrained systems)
\item explicit separation between ``original'' physics and ``differentiable'' variants
\end{itemize}

The refactor adds a substantial set of \texttt{*\_diff} routines and numerical utilities
(smoothers, clamps, overshoot fixes, implicit solve scaffolding).

\subsubsection{Why this may change results}

These changes are designed to improve stability and derivative behavior, but they can also
alter trajectories:
\begin{itemize}[leftmargin=*]
\item Clamps and overshoot protection modify fluxes/states when constraints are violated.
\item Smoothers modify transitions that were previously sharp/discontinuous.
\item Implicit solve tolerances and convergence behavior can differ from explicit stepping.
\end{itemize}

Therefore, when strict bitwise regression is required, it is useful to ensure these features
can be toggled or to run equivalent pathways for comparison.

\subsection{Output Hygiene: Reducing Ad Hoc Print Statements}

\subsubsection{Why pervasive print/debug statements were a problem}

The legacy code contained many ad hoc \texttt{print} (and related) statements inserted over
years of development. While individually helpful during debugging, collectively they became a
maintenance and operational burden:

\begin{itemize}
\item \textbf{Signal-to-noise degradation:} important warnings and true error messages were
      buried in routine diagnostic output.
\item \textbf{Non-deterministic verbosity:} output depended on the exact execution pathway
      and developer preferences, complicating regression comparisons.
\item \textbf{Difficult automation:} noisy logs make it harder to detect failure patterns
      in batch runs or CI-style workflows.
\end{itemize}

\subsubsection{Why the refactor is better}

Removing large volumes of ad hoc output is not merely cosmetic: it improves debuggability by
making the remaining output more meaningful. The refactor moves toward an implicit contract:

\begin{itemize}
\item \textbf{Routine execution is quiet by default.}
\item \textbf{Errors and validation failures are explicit and actionable.}
\item \textbf{Diagnostics are intentional (ideally controlled by flags or structured logging).}
\end{itemize}

This makes calibration runs, ensemble runs, and regression testing substantially easier to manage.

% =====================================================
\section{Migration Roadmap (toward the end-state design)}

\subsection{Status}

The redesign is already in place for the driver/evaluation split and for explicit time handling.
However, some legacy routines continue to rely on global modules. The compatibility layer in
\texttt{share/*\_data.f90} exists to keep these routines operational while migration continues.

\subsection{Planned milestones}

\begin{enumerate}[leftmargin=*]
\item \textbf{Complete container migration:} update remaining routines to accept \texttt{info/work/domain}
      (or narrow, purpose-built structures) instead of reading from globals.
\item \textbf{Reduce compatibility surface:} progressively shrink what is exported to \texttt{share/*\_data}
      modules, and add warnings or comments indicating ``do not use in new code''.
\item \textbf{Remove compatibility modules:} once call sites are migrated, remove the legacy global storage
      modules entirely.
\item \textbf{Formalize testing:} establish a regression test suite that checks time indexing, forcing alignment,
      basic hydrograph behavior, and metric reproducibility across run modes.
\item \textbf{Document extension points:} add developer docs describing how to add forcing variables, new metrics,
      and new physics modules without reintroducing global coupling.
\end{enumerate}

\subsection{Suggested regression checklist}

\begin{enumerate}[leftmargin=*]
\item Validate \texttt{info\%time}: units, reference date, first/last timesteps, \texttt{deltim\_days}.
\item Validate window indices: \texttt{sim\_beg/sim\_end} and \texttt{eval\_beg/eval\_end}.
\item Validate forcing alignment: sample a few variables at a few timesteps and confirm expected values.
\item Validate deterministic evaluation: compare outlet discharge and key states to a known baseline.
\item Validate calibration callback: ensure SCE evaluations are consistent and do not leak state between calls.
\end{enumerate}

% =====================================================
\section{Conclusion}

This development represents both a structural re-platforming and a major capability
expansion of the FUSE modeling framework. While the architectural refactor replaces
an implicit global-state design with structured run-state containers, modular
driver/evaluation logic, and centralized time and domain handling, these changes
were undertaken to enable substantially new functionality.

Most notably, the framework now supports a fully differentiable physics pathway
with analytical Jacobians propagated through snow and soil processes, and a
purpose-built fixed-step implicit solver tailored for deterministic,
derivative-consistent execution. In addition, the introduction of a generalized
spatial-domain abstraction enables unified treatment of lumped, distributed,
and gridded configurations, along with structured spatial decomposition suitable
for scalable parallel execution.

The legacy architecture, while historically productive and scientifically
influential, relied on implicit shared state and convention-based execution.
Under modern requirements—gradient-based calibration, reproducible execution,
and extensible spatial deployment—these patterns became limiting. The redesigned
architecture addresses hidden dependencies, initialization-order fragility, and
high extension cost by making state explicit, separating orchestration from
process execution, and clarifying data flow throughout the call graph.

The compatibility layer (\texttt{share/*\_data}) is intentionally transitional and
should be removed once migration is complete. The resulting end-state is not
merely cleaner code; it is a modeling framework capable of differentiable
execution, robust fixed-step integration, flexible spatial deployment, and
future integration with adjoint-based and machine-learning methodologies.

In this sense, the refactor should be viewed not as code cleanup, but as an
evolution of FUSE from a flexible forward-simulation platform into a structured,
extensible system designed for the next generation of hydrologic modeling
research.

\end{document}
